\section{Sums, Temperature and the Puzzles Revisited}\label{sec:sums}

Most endgames are sums of components. Two principles guide play in them. By \emph{number avoidance}, in a sum of
numbers and games that are not numbers, a player who can win can do so without moving in a
number~\cite{WinningWays,Siegel2013}. By the heuristic of the \emph{hottest move}, one moves in the component of
highest temperature: for a switch $\game{x}{y}$ between numbers $x>y$, the temperature $(x-y)/2$ is what each
side gains by moving there first.

\subsection{The Puzzles Revisited}

Table~\ref{tab:puzzles} compares the values first claimed for three puzzles of~\cite{Study2025} with the exact
ones: the earlier walk-throughs pick the right winners, but not always the right values. In Puzzle~1, black's
h6$\times$a5 across the seam leaves $\star$, so the h/a component is hot, not the number $1$. The
pawns of Puzzle~2 form an infinitesimal with mean $0$, and its king group (Fig.~\ref{fig:kinggroup72}) is a
zugzwang for black only, on the ordinary board \lab{C2} (our solver and the Lichess tablebase) and on the
cylinder \lab{C1}, so the real result of this ten-man position is open. Puzzle~3 is illegal as drawn; on files c and d its value is $\tinyg{1}$, and the legal copy
one rank up, \pos{W c2 c4 d4 / B c3 d5 d6 d7}, has the same value (Section~S3 shows the puzzles).

\begin{table}[!t]
\centering
\caption{Puzzles of the 2025 Study: Claimed and Exact Values \lab{C2}}
\label{tab:puzzles}
\begin{tabular}{@{}lll@{}}
\toprule
Part & Claimed~\cite{Study2025} & Exact \\
\midrule
1: c-file & $\pm1$ & $\pm1$ \\
1: f-file & $\game{\game{2}{1}}{0}$ & the same (mean $3/4$) \\
1: files h, a & $1$ & $\game{+\BIG}{\star}$ (hot) \\
1: total & $1\pm1+\game{\game{2}{1}}{0}$ & white wins either way \\
2: pawns & $\pm1$ & infinitesimal, mean $0$ \\
2: kings & $0$ & not mutual \\
3: files c, d & $1$ & $\tinyg{1}$ \\
3: total & $2$ & $+\BIG$ before any move \\
\bottomrule
\end{tabular}
\end{table}

\subsection{When Is the Sum Model Right?}\label{sec:sums-hss}

In chess the first passed pawn ends the game; in a sum, a component with a passed pawn becomes $\pm\BIG$ and
play goes on elsewhere, so a $+\BIG$ can cancel a $-\BIG$. In $\game{\BIG}{0}+\game{0}{-\BIG}$ the first player
makes a passed pawn and wins the chess game, but the sum is $0$ for every $\BIG$, since the second player
answers in the other component \lab{P}. We write $W(X)$ for the whole-board game of a pawn position $X$, which
ends at the first passed pawn, and $S(X)$ for its sum model, which after every move re-splits the pawns into
components along empty files (one ring if no file is empty), replaces a component in which a side has a passed
pawn by $\pm\BIG$ (by the race rule if both have one) and adds the components up. We call $W(X)$ a $\nim{k}$ when
$W(X)+\nim{h}$ has the outcome of $\nim{k}+\nim{h}$ for every heap $h$; our scans and the wall proofs test
exactly this.

\begin{definition}[Hereditary sum safety]\label{def:hss}
A component carries a \emph{threat} for a side if a position reachable from it gives that side a passed pawn. A
position $X$ is \emph{hereditarily sum-safe} (HSS) if no position reachable from $X$ by moves of either side,
in any order, before the first passed pawn, $X$ included, has two components with threats for opposite sides.
\end{definition}

\begin{theorem}[Sum safety; \lab{S}]\label{thm:hss}
Suppose that $X$ is HSS, that no position reachable from $X$ with two or more components has a usable en
passant right (condition (E)), and that $\BIG$ exceeds the number of moves either side can make. Then $S(X)$ and
$W(X)$ have the same outcome.
\end{theorem}

\begin{proof}[Proof sketch]
The models differ only after a passed pawn appears, say for white. By HSS no other component can then give black
one, so black runs out of moves before white's $+\BIG$ does. The full proof, by induction with (E), is in
Section~S12 of the supplementary material.
\end{proof}

In a random sample of $1{,}500$ positions, all $1{,}297$ HSS ones agree \lab{C1}, and with $\BIG=256$ the move
bound always holds on boards up to $8\times8$. HSS cannot be dropped (Proposition~\ref{prop:sumcounter}); without
(E) the decomposition already fails in the pure pawn game, without the passed-pawn rule \lab{C1} (Section~S3).

\begin{proposition}[\lab{C2}]\label{prop:sumcounter}
A test of the top-level components alone is not enough: on $\cyl{8}$ with en passant off, $X=$\,\pos{W b2 d4 g3
h5 / B a7 c7 f7 g4} is a single component, spanning files f to d across the seam, so such a test passes it; yet
$S(X)$ is a black win whoever starts and $W(X)$ a first-player win. Moreover $X$ is not HSS: components with
opposite threats appear later in its tree, so Theorem~\ref{thm:hss} does not apply.
\end{proposition}

\emph{Two models.} Our CGT engine computes $S$, as Elkies does~\cite{Elkies1996}, and the independent
whole-board solver $W$. They agree on every headline value: the $1/8$ and $\upsec\star$ in all nine settings,
and the seven-pawn $\nim{2}$ in the default setting. Elsewhere they can differ, even in outcome (Section~S3);
there we name the model we mean.
% SLOT: real-king sum test
